\section{Integral Theorems and Methods}
\subsection{Linearity, or sum rule, difference rule, and constant multiple rule}
If functions $f$ and $g$ are integrable over an interval $I$ and $a,b\in\mathbb{R}$, then $af(x)+bg(x)$ is also integrable over $I$ and its integral is given by
\[\int_Iaf(x)+bg(x)\dd{x}=a\int_If(x)\dd{x}+b\int_Ig(x)\dd{x}.\]
\subsection{Integrals of Symmetric Functions}
Suppose real function $f$ is Riemann integrable on $[-a,a]$.
\bit
\item If $f$ is an even function, that is, $f(-x)=f(x)$, then $\int_{-a}^af(x)\dd{x}=2\int_0^af(x)\dd{x}$.
\item If $f$ is an of function, that is, $f(-x)=-f(x)$, then $\int_{-a}^af(x)\dd{x}=0$.
\eit
\subsection{Mean Value Theorem (MVT) for Integrals}
If a function $f\colon D\subseteq\bbR\to\bbR$ is continuous on a closed interval $[a,b]$ with $a<b$, then
\[\exists c\in (a, b)\text{\ s.t.\ }f(c)=\frac{1}{b-a}\int_a^bf(x)\dd{x}.\]
\begin{proof}
By EVT, $f$ must have maximum and minimum on $[a,b]$.

Let
\[M=\max_{x\in [a,b]}f(x),\quad m=\min_{x\in [a,b]}f(x).\]
Then for any $x\in[a,b]$:
\[m\leq f(x)\leq M.\]
Integrate it over $[a,b]$:
\[m(b-a)=\int_a^bm\dd{x}\leq\int_a^bf(x)\dd{x}\leq\int_a^bM\dd{x}=M(b-a).\]
\[m\leq\frac{1}{b-a}\int_a^bf(x)\dd{x}\leq M.\]
Since $f$ is continuous, by IVT, there exists $c\in[a,b]$ such that
\[f(c)=\frac{1}{b-a}\int_a^bf(x)\dd{x}.\]
\end{proof}
\subsection{Fundamental Theorem of Calculus (FTC) (微積分基本定理)}
\subsubsection{Fundamental Theorem of Calculus Part First or First Fundamental Theorem of Calculus}
Let function $f\colon D\subseteq\bbR\to\bbR$ be continuous on a closed interval $[a,b]$, and $F$ be the function defined, for all $x\in[a,b]$ by
\[F(x)=\int_a^xf(t)\dd{t}.\]
Then $F$ Lipschitz continuous on $[a,b]$ and differentiable on $(a,b)$ and
\[F'(x)=f(x)\]
for all $x\in (a,b)$.
\begin{proof}
Lipschitz continuity: Let $(x,y)\in(a,b)$ with $x<y$. Let
\[M=\max_{t\in[a,b]}\abs{f(t)}.\]
Then,
\[F(y)-F(x)=\int_a^yf(t)\dd{t}-\int_a^xf(t)\dd{t}=\int_x^yf(t)\dd{t}.\]
\[\abs{F(y)+F(x)}=\abs{\int_x^yf(t)\dd{t}}\leq\int_x^y\abs{f(t)}\dd{t}\leq\int_x^yM\dd{t}=M(y-x).\]
Thus $F$ is Lipschitz continuous with constant $M$.

Differentiability:
\[F'(x)=\lim_{h\to 0}\frac{\int_a^{x+h}f(t)\dd{t}-\int_a^xf(t)\dd{t}}{h}=\lim_{h\to 0}\frac{\int_x^{x+h}f(t)\dd{t}}{h}.\]
By MVT for integrals, there exists $c_h\in[x,x+h]$ (or $[x+h,x]$ for h<0) such that
\[f(c_h)=\frac{\int_x^{x+h}f(t)\dd{t}}{h}.\]
By continuity of $f$,
\[\lim_{h\to 0}c_h=x.\]
Hence,
\[\lim_{h\to 0}\frac{\int_x^{x+h}f(t)\dd{t}}{h}=f(x).\]
\end{proof}
\subsubsection{Fundamental Theorem of Calculus Part Second, Second Fundamental Theorem of Calculus, or Newton–Leibniz Theorem}
Let function $f\colon I\subseteq\bbR\to\bbR$ be Riemann integrable on a closed interval $[a,b]$, and $F\colon J\subseteq\bbR\to\bbR$ be a function continuous on $[a,b]$ and differentiable on $(a,b)$ such that
\[F'(x)=f(x)\]
for all $x\in (a,b)$.

Then
\[\int_a^bf(x)\dd{x}=F(b)-F(a).\]
\begin{proof}
Let $P(x,n)$ be a partition of $[a,b]$. By MVT, for each subinterval $[x_i,x_{i+1}]$, there exists $c_i\in[x_i,x_{i+1}]$ such that
\[F(x_{i+1})-F(x_i)=F'(c_i)(x_{i+1}-x_i)=f(c_i)(x_{i+1}-x_i).\]
So for any $n\in\bbN$:
\[\sum_{i=0}^{n-1}f(c_i)(x_{i+1}-x_i)=\sum_{i=0}^{n-1}(F(x_{i+1})-F(x_i))=F(x_n)-F(x_0)=F(b)-F(a).\]
Since $f$ is Riemann integrable
\[\int_a^bf(x)\dd{x}=\lim_{n\to\infty}\sum_{i=0}^{n-1}f(c_i)(x_{i+1}-x_i)\]
exists.
\end{proof}
\subsection{Integrability}
\subsubsection{Lebesgue-Vitali theorem (of characterization of the Riemann integrable functions)}
A function is Riemann integrable (i.e. Darboux integrable) over a closed interval $I$ iff it is bounded on $I$ and continuous a.e. in $I$.
\begin{proof}
PLACEHOLDER
\end{proof}
\subsubsection{Riemann–Stieltjes integrability}
The Riemann–Stieltjes integral
\[\int_{x=a}^bf(x)\dd{g(x)}\]
exists if $g$ is of a bounded variation on $[a,b]$, $f$ is bounded on $[a,b]$, $f$ is continuous a.e. on $[a,b]$, and there doesn't exist $x\in[a,b]$ such that $f$ and $g$ both have irremovable discontinuity at $x$.
\begin{proof}
PLACEHOLDER
\end{proof}
\subsubsection{PLACEHOLDER}
PLACEHOLDER
\ssc{Absolutely continuous antiderivative}
If function $F(x)$ is absolutely continuous on an open interval $(a,b)$, then for any $f_1,f_2\in L^1((a,b))$ such that
\[f(a)+\int_a^xf_1(t)\dd{t}=\int_a^xf_2(t)\dd{t}=F(x),\]
we have
\[f_1(x)=f_2(x)\]
a.e. on $(a,b)$, and
\[f_1(x)=F'(x)\]
a.e. on $(a,b)$. (Note that $F$ is differentiable a.e. since it's absolutely continuous.)
\subsection{Integration by substitution (代換積分法), integration by change of variables (換元積分法), u-substitution (u-代換), reverse chain rule, substitution theroem (代換定理), change of variables theorem (換元定理), or transformation theorem (變換定理)}
\subsubsection{Univariate form}
Let $g:\,I\subseteq\mathbb{R}\to\mathbb{R}$ be injective and differentiable on $[a,b]\subseteq I$, with $g'$ being integrable on $[a,b]$, and $f:\,K\supseteq g([a,b])\to\mathbb{R}$ be integrable on $g([a,b])$. Then:
\[\int_a^bf\circ g(x)\cdot g'(x)\dd{x}=\int_{g(a)}^{g(b)}f(u)\,\mathrm{d}u,\]
and
\[\int f\circ g(x)\cdot g'(x)\dd{x}=\int f(u)\,\mathrm{d}u.\]
Substitutions from RHS to LHS are also called inverse substitution.
\begin{proof}
Consider the interval \([a, b]\) partitioned as
\[P = \{a = x_0 < x_1 < \dots  < x_n = b\}.\]
For each subinterval \([x_{i-1}, x_i]\), let \(\xi_i \in [x_{i-1}, x_i]\) be a sample point. The Riemann sum for the left-hand side of the integral is:
\[S_P = \sum_{i=1}^n f(g(\xi_i)) g'(\xi_i) (x_i - x_{i-1}).\]
Since \(g\) is injective and differentiable, it is either strictly increasing or strictly decreasing on \([a, b]\). Assume \(g\) is strictly increasing (the proof for \(g\) strictly decreasing follows similarly).

Under this assumption, \(g\) maps \([a, b]\) to \([g(a), g(b)]\) with \(g(a) < g(b)\). Let \([g(a), g(b)]\) be partitioned as
\[Q = \{g(a) = u_0 < u_1 < \dots  < u_m = g(b)\},\]
where \(u_i = g(x_i)\).

The Riemann sum for the right-hand side of the integral is:
\[T_Q = \sum_{i=1}^n f(u_i) (u_i - u_{i-1}).\]
Since \(u_i = g(x_i)\) and \(g'(\xi_i) \approx \frac{g(x_i) - g(x_{i-1})}{x_i - x_{i-1}}\), we rewrite:
\[ g'(\xi_i) (x_i - x_{i-1}) \approx g(x_i) - g(x_{i-1}) = u_i - u_{i-1}. \]
Thus, the left-hand Riemann sum \(S_P\) becomes:
\[ S_P = \sum_{i=1}^n f(g(\xi_i)) (u_i - u_{i-1}),\]
which matches the structure of the right-hand Riemann sum \(T_Q\) if we let \(\xi_i = g^{-1}(u_i)\).

As the partition \(P\) of \([a, b]\) gets finer, the Riemann sum \(S_P\) converges to \(\int_a^b f(g(x)) g'(x) \, \mathrm{d}x\). Similarly, as the partition \(Q\) of \([g(a), g(b)]\) gets finer, the Riemann sum \(T_Q\) converges to \(\int_{g(a)}^{g(b)} f(u) \, \mathrm{d}u\).
\end{proof}
\subsubsection{Multivariate form}
Let $\mathbf{T}\colon I\subseteq\mathbb{R}^n\to\mathbb{R}^n;\;\mb{x}\mapsto\mb{y}$ be injective and differentiable on $D\subseteq I$, its Jacobian matrix $\nabla\mathbf{T}$ being continuous on $D$, and $f\colon K\supseteq\mathbf{T}(D)\to\mathbb{R}$ be integrable on $\mathbf{T}(D)$. Then
\[\int_Df\qty(\mb{x})\abs{\det\left(\nabla\mathbf{T}\right)}\dd{V}=\int_{\mathbf{T}(D)}f\qty(\mb{T}^{-1}(\mb{y}))\dd{V}\]
and
\[\int f\qty(\mb{x})\abs{\det\left(\nabla\mathbf{T}\right)}\dd{V}=\int f\qty(\mb{T}^{-1}(\mb{y}))\dd{V}.\]
Substitutions from RHS to LHS are also called inverse substitution.
\bpr
PLACEHOLDER
\epr
If $T^{-1}$ is differentiable,
\[\det\left(\nabla\mathbf{T}\right)\det\left(\nabla(\mathbf{T}^{-1})\right)=1\]
Thus,
\[\int_Df\qty(\mb{x})\dd{V}=\int_{\mathbf{T}(D)}f\qty(\mb{T}^{-1}(\mb{y}))\abs{\det\left(\nabla(\mathbf{T}^{-1})\right)}\dd{V}\]
and
\[\int f\qty(\mb{x})l\dd{V}=\int f\qty(\mb{T}^{-1}(\mb{y}))\abs{\det\left(\nabla(\mathbf{T}^{-1})\right)}\dd{V}.\]
\subsubsection{Measure theory form}
Let $X$ be a locally compact Hausdorff space equipped with a finite Radon measure $μ$, and let $Y$ be a \text{\textsigma}-compact Hausdorff space with a \text{\textsigma}-finite Radon measure $\rho$. Let $\phi\colon X\to Y$ be an absolutely continuous function. Then there exists a real-valued Borel measurable function $w$ on $X$ such that, for every function $f\colon Y\to\mathbb{R}$ in $L^1(Y)$, the function $(f\circ\phi)\cdot w$ is in $L^1(X)$, and for every open subset $U$ of $X$
\[\int_U(f\circ\phi)(x)\cdot w(x)\dd{\mu(x)}=\int_{\phi(U)}f(y)\dd{\rho(y)}.\]
Substitutions from RHS to LHS are also called inverse substitution.

Furthermore, there exists some real-valued Borel measurable function $g$ such that 
\[w(x)=(g\circ\phi)(x).\]
\bpr
PLACEHOLDER
\epr
\subsection{Integration by parts (IBP) (分部積分法) or partial integration (部分積分法)}
\subsubsection{Theorem}
\[\frac{\mathrm{d}}{\mathrm{d}x}\prod_{i=1}^nf_i(x)=\sum_{j=1}^n\left(\frac{\mathrm{d}f_j(x)}{\mathrm{d}x}\frac{\prod_{\substack{i=1\\i\neq j}}^n f_i(x)}{f_j(x)}\right)\]
For example,
\[\int_{\Omega} u\dd{v} = \qty(u v)\big\vert_{\Omega} - \int_{\Omega} v\dd{u}.\]
\subsubsection{Application}
Integration by parts is a heuristic rather than a purely mechanical process for solving integrals; given a single function to integrate, the typical strategy is to carefully separate this single function into a product of two functions such that the residual integral from the integration by parts formula is easier to evaluate than the single function.

The DETAIL rule or the LIATE rule is a rule of thumb for integration by parts. It involves choosing as u the function that comes first in the following list:
\begin{itemize}
\item L: Logarithmic function
\item I: Inverse trigonometric function
\item A: Algebraic function
\item T: Trigonometric function
\item E: Exponential function
\end{itemize}
\subsubsection{Notation for two parts}
We write the integration of
\[\int f(x)\dd{x}\]
by parts as:
\ben
\item Let
\[u=g(x),\quad\dd{v}=h(x)\dd{x}\]
with some functions $g(x)$ and $h(x)$ such that
\[f(x)\dd{x}=u\dd{v}.\]
\item Compute
\[\dd{u}=g'(x)\dd{x},\quad v=\int h(x)\dd{x}\]
with the constant of integration of $\int h(x)\dd{x}$ arbitrarily chosen to be assigned to $v$.
\item Integral by parts
\[\int f(x)\dd{x}=uv-\int v\dd{u}.\]
\item If $\int v\dd{u}$ is simplified as
\[\int v\dd{u}=k\int f(x)\dd{x}+I(x)+C\]
with some constant $k\notin\{-1,0\}$, we write
\[\ba
\int f(x)\dd{x}&=uv-\int v\dd{u}\\
&=uv-k\int f(x)\dd{x}-I(x)+C\\
&=\frac{1}{k+1}\qty(uv-I(x))+C.
\ea\]
\een
Take
\[\int e^{ax}\sin(bx)\dd{x}\]
for example.

Let:
\[u=\sin(bx),\quad\dd{v}=e^{ax}\dd{x}.\]
\[\dd{u}=b\cos(bx)\dd{x},\quad v=\int e^{ax}\dd{x}=\frac{1}{a}e^{ax}.\]
\[\ba
\int e^{ax}\sin(bx)\dd{x}&=\int u\dd{v}=uv-\int v\dd{u}\\
&=\frac{1}{a}e^{ax}\sin(bx)-\frac{b}{a}\int e^{ax}\cos(bx)\dd{x}.
\ea\]
\[w=\cos(bx).\]
\[\dd{w}=-b\sin(bx).\]
\[\ba
\int e^{ax}\cos(bx)\dd{x}&=\int w\dd{v}=wv-\int v\dd{w}\\
&=\frac{1}{a}e^{ax}\cos(bx)+\frac{b}{a}\int e^{ax}\sin(bx)\dd{x}.
\ea\]
\[\ba
\int e^{ax}\sin(bx)\dd{x}&=\frac{1}{a}e^{ax}\sin(bx)-\frac{b}{a}\qty(\frac{1}{a}e^{ax}\cos(bx)+\frac{b}{a}\int e^{ax}\sin(bx)\dd{x})\\
&=\frac{1}{a}e^{ax}\sin(bx)-\frac{b}{a^2}e^{ax}\cos(bx)-\frac{b^2}{a^2}\int e^{ax}\sin(bx)\dd{x}\\
&=\frac{a^2}{a^2+b^2}\qty(\frac{1}{a}e^{ax}\sin(bx)-\frac{b}{a^2}e^{ax}\cos(bx))+C\\
&=\frac{e^{ax}}{a^2+b^2}\qty(a\sin(bx)-b\cos(bx))+C.
\ea\]
\subsection{Riemann–Stieltjes integral theorems}
\subsubsection{Integration by parts (IBP) or partial integration of Riemann–Stieltjes integral}
\[\int_{x=a}^bf(x)\dd{g(x)}=f(b)g(b)-f(a)g(a)-\int_{x=a}^bg(x)\dd{f(x)}.\]
\begin{proof}
PLACEHOLDER
\end{proof}
\subsubsection{Riemann–Stieltjes integral with respect to differential function}
If $f(x)$ is bounded on $[a,b]$, $g(x)$ is monotonic on $[a,b]$, and $g'(x)$ exists and is Riemann integrable on $[a,b]$, then the Riemann–Stieltjes integral is related to the Riemann integral by
\[\int_{x=a}^bf(x)\dd{g(x)}=\int_a^bf(x)g'(x)\dd{x}.\]
\begin{proof}
PLACEHOLDER
\end{proof}
\subsubsection{Riemann–Stieltjes integral with respect to step function}
If $a<s<b$,
\[g(x)=\bcs 0,\quad&a\leq x\leq s\\
1,\quad&b\geq x>s\ecs,\]
and $f\colon X\supseteq[a,b]\to\bbR$ is continuous at $s$, then
\[\int_{x=a}^bf(x)\dd{g(x)}=f(s).\]
\begin{proof}
PLACEHOLDER
\end{proof}
\subsection{General Method of Integration of Rational Functions}
\subsubsection{$\frac{A}{x-a}$}
\[\int\frac{A}{x-a}\dd{x}=A\ln|x-a|+C,\quad a,A\in\bbR\]
\subsubsection{$\frac{A}{\qty(x-a)^g}$}
\[\int\frac{A}{\qty(x-a)^g}\dd{x}=-\frac{A}{(g-1)\qty(x-a)^{g-1}}+C,\quad a,A\in\bbR\land g\in\bbN_{>1}\]
\subsubsection{$\frac{Ax+B}{x^2+bx+c}$}
\[\int\frac{Ax+B}{x^2+bx+c}=\frac{A}{2}\ln|x^2+bx+c|+\frac{2B-bA}{\sqrt{4c-b^2}}\arctan(\frac{2x+b}{\sqrt{4c-b^2}})+C,\quad b,c,A,B\in\bbR\land b^2<4c\]
\begin{proof}
\[\ba
\int\frac{Ax+B}{x^2+bx+c}&=\int\frac{A}{2}\frac{2x+b}{x^2+bx+c}+\frac{B-\frac{bA}{2}}{\qty(x+\frac{b}{2})^2+\qty(c-\frac{b^2}{4})}\dd{x}\\
&=\frac{A}{2}\ln|x^2+bx+c|+\qty(B-\frac{bA}{2})\int\frac{1}{u^2+\qty(c-\frac{b^2}{4})}\dd{u},\quad u=x+\frac{b}{2}\\
&=\frac{A}{2}\ln|x^2+bx+c|+\frac{2B-bA}{\sqrt{4c-b^2}}\arctan(\frac{2x+b}{\sqrt{4c-b^2}})+C
\ea\]
\subsubsection{$\frac{1}{\qty(x^2+bx+c)^h}$}
\[\int\frac{1}{\qty(x^2+bx+c)^h}\dd{x}=\frac{1}{\qty(4c-b^2)(h-1)}\frac{2x+b}{\qty(x^2+bx+c)^{h-1}}+\frac{4h-6}{\qty(4c-b^2)(h-1)}\int\frac{1}{\qty(x^2+bx+c)^{h-1}}\dd{x},\quad b,c\in\bbR\land b^2<4c\land h\in\bbN_{>1}\]
\begin{proof}
\[u=x+\frac{b}{2},\quad\alpha=\sqrt{c-\frac{b^2}{4}}\]
\[\ba 
\int\frac{1}{\qty(x^2+bx+c)^h}\dd{x}&=\int\frac{1}{\qty(u^2+\alpha^2)^h}\dd{u}\\
&=\frac{1}{\alpha^2}\int\frac{u^2+\alpha^2-u^2}{\qty(u^2+\alpha^2)^h}\dd{u}\\
&=\frac{1}{\alpha^2}\int\frac{1}{\qty(u^2+\alpha^2)^{h-1}}\dd{u}-\frac{1}{\alpha^2}\int\frac{u^2}{\qty(u^2+\alpha^2)^h}\dd{u}
\ea\]
\[\ba
\dv{}{u}\qty(\frac{u}{\qty(u^2+\alpha^2)^{h-1}})&=\frac{\qty(u^2+\alpha^2)^{h-1}-u(h-1)\qty(u^2+\alpha^2)^{h-2}(2u)}{\qty(u^2+\alpha^2)^{2(h-1)}}\\
&=\frac{1}{\qty(u^2+\alpha^2)^{h-1}}-
2(h-1)\frac{u^2}{\qty(u^2+\alpha^2)^h}
\ea\]
\[\frac{u^2}{\qty(u^2+\alpha^2)^h}=\frac{1}{2(h-1)}\qty(\frac{1}{\qty(u^2+\alpha^2)^{h-1}}-\dv{}{u}\qty(\frac{u}{\qty(u^2+\alpha^2)^{h-1}}))\]
\[\int\frac{u^2}{\qty(u^2+\alpha^2)^h}\dd{u}=\frac{1}{2(h-1)}\qty(\int\frac{1}{\qty(u^2+\alpha^2)^{h-1}}\dd{u}-\frac{u}{\qty(u^2+\alpha^2)^{h-1}})\]
\[\ba  
\int\frac{1}{\qty(x^2+bx+c)^h}\dd{x}&=\frac{1}{\alpha^2}\int\frac{1}{\qty(u^2+\alpha^2)^{h-1}}\dd{u}-\frac{1}{\alpha^2}\int\frac{u^2}{\qty(u^2+\alpha^2)^h}\dd{u}\\
&=\frac{1}{\alpha^2}\int\frac{1}{\qty(u^2+\alpha^2)^{h-1}}\dd{u}-\frac{1}{2\alpha^2(h-1)}\qty(\int\frac{1}{\qty(u^2+\alpha^2)^{h-1}}\dd{u}-\frac{u}{\qty(u^2+\alpha^2)^{h-1}})\\
&=\frac{2h-3}{2\alpha^2(h-1)}\int\frac{1}{\qty(u^2+\alpha^2)^{h-1}}\dd{u}+\frac{1}{2\alpha^2(h-1)}\frac{u}{\qty(u^2+\alpha^2)^{h-1}}\\
&=\frac{4h-6}{\qty(4c-b^2)(h-1)}\int\frac{1}{\qty(x^2+bx+c)^{h-1}}\dd{x}+\frac{1}{\qty(4c-b^2)(h-1)}\frac{2x+b}{\qty(x^2+bx+c)^{h-1}}
\ea\]
\end{proof}
\subsubsection{$\frac{Ax+B}{\qty(x^2+bx+c)^h}$}
\[\int\frac{Ax+B}{\qty(x^2+bx+c)^h}\dd{x}=-\frac{A}{2(h-1)\qty(x^2+bx+c)^{h-1}}+\qty(B-\frac{bA}{2})\int\frac{1}{\qty(x^2+bx+c)^h}\dd{x},\quad b,c,A,B\in\bbR\land b^2<4c\land h\in\bbN_{>1}\]
\subsubsection{Integration of Rational Functions by Partial Fractions}
To compute the integral of any arbitrary rational function $f(x)$,
\ben
\item Express it as
\[f(x)=S(x)+\frac{P(x)}{Q(x)}\]
where $S(x),P(x),Q(x)$ are polynomials and $\deg(P)<\deg(Q)$.
\item Find $q,a_j,b_k,c_k\in\bbR$, $m,n\in\bbN_0$, $r_j,s_k\in\bbN$ and $b_k^{\pht{k}2}-4c_k<0$ for all $k$ such that
\[Q(x)=q\prod_{j=1}^m\qty(x-a_j)^{r_j}\prod_{k=1}^n\qty(x^2+b_kx+c_k)^{s_k}.\]
Note that all polynomials can be expressed as a product of linear factors and irreducible quadratic factors.
\item Express $\frac{P(x)}{Q(x)}$ as
\[\frac{P(x)}{Q(x)}=p\qty(\sum_{j=1}^m\sum_{g=1}^{r_j}\frac{A_{j_g}}{\qty(x-a_j)^g}+\sum_{k=1}^n\sum_{h=1}^{s_k}\frac{B_{k_h}x+C_{k_h}}{\qty(x^2+b_kx+c_k)^h})\]
where $p,A_{j_g},B_{k_h},C_{k_h}\in\bbR$.
\item Integrate $S(x)$ and each term of $\frac{A_{j_g}}{\qty(x-a_j)^g}$ and $\frac{B_{k_h}x+C_{k_h}}{\qty(x^2+b_kx+c_k)^h}$ respectively.
\een
\subsection{Substitution for irrational function integrand}
\subsubsection{Direct substitution}
When computing an integral of
\[R(x,\sqrt[n]{g(x)})\]
with respect to $x$ where $R$ and $g$ are rational functions, the substitution $u=\sqrt[n]{g(x)}$ may be effective.
\subsubsection{Square root of quadratic polynomial}
When computing an integral of
\[R\qty(x,\sqrt{ax^2+bx+c})\]
with respect to $x$ where $R$ is a rational function and $a\neq 0$, it may be effective to
\ben
\item let $h=\frac{b}{2|a|}$ and $k=-\frac{b^2}{4a^2}+\frac{c}{a}$, rewrite $\sqrt{ax^2+bx+c}$ as
\[\sqrt{ax^2+bx+c}=\sqrt{|ak|}\sqrt{\sgn(a)\qty(\frac{x}{\sqrt{|k|}}+\frac{h}{\sqrt{|k|}})^2+\sgn(k)},\]
\iten use substitution $u=\frac{x}{\sqrt{|k|}}+\frac{h}{\sqrt{|k|}}$, which makes the integrand,
\[\sqrt{|k|}R\qty(\sqrt{|k|}u-h,\sqrt{|ak|}\sqrt{\sgn(a)u^2+\sgn(k)}),\]
which is a rational function of $u$ and $\sqrt{\sgn(a)u^2+\sgn(k)}$, and,
\item depending on $\sgn(a)$ and $\sgn(k)$,
\bit
\item $k=0$: the integrand is a rational function of $u$,
\item $a>0$ and $k>0$: try the substitution $u=\tan(t)$, $u=\cot(t)$, $u=\sinh(t)$, or $u=\csch(t)$,
\item $a>0$ and $k<0$: try the substitution $u=\sec(t)$, $u=\csc(t)$, $u=\cosh(t)$, or $u=\coth(t)$,
\item $a<0$ and $k>0$: try the substitution $u=\sin(t)$, $u=\cos(t)$, $u=\tanh(t)$, or $u=\sech(t)$.
\eit
\een
\subsubsection{Euler substitutions}
When computing an integral of
\[R(x,\sqrt{ax^2+bx+c})\]
with respect to $x$ where $R$ is a rational function, the three Euler('s) substitutions may be effective.
\ben
\item First substitution: For $a\geq 0$, substitute
\[\sqrt{ax^2+bx+c}=\pm x\sqrt{a}+t\]
\[x=\frac{c-t^2}{\pm 2t\sqrt{a}-b}\]
where either the positive sign or the negative sign can be chosen.
\item Second substitution: For $c\geq 0$, substitute
\[\sqrt{ax^2+bx+c}=xt\pm\sqrt{c}\]
\[x=\pm 2t\sqrt{c}-b}\]
where either the positive sign or the negative sign can be chosen.
\item Third substitution: For $b^2-4ac\geq 0$, let
\[ax^2+bx+c=a(x-\alpha)(x-\beta),\]
substitute
\[\sqrt{ax^2+bx+c}=(x-\alpha)t\]
\[x=\frac{\alpha\beta-\alpha t^2}{a-t^2}\]
\een
\subsection{Function of trigonometric functions integral}
\subsubsection{Tangent half-angle substitution, Weierstrass substitution, or universal trigonometric substitution}
When computing an integral of
\[f(\sin(x),\cos(x))\]
with respect to $x$, the substitution
\[t=\tan(\frac{x}{2})\]
gives
\[\int f(\sin(x),\cos(x))=\int f\qty(\frac{2t}{1+t^2},\frac{1-t^2}{1+t^2})\frac{2}{1+t^2}\dd{t}.\]
\subsubsection{Product of sine and cosine functions}
\[\int\sin^m(x)\cos^n(x)\dd{x},\quad m,n\in\bbN.\]
\begin{itemize}
\item For odd $n$, use $\cos^2(x)=1-\sin^2(x)$ and then substitute $u=\sin(x),\quad\dd{u}=\cos(x)\dd{x}$.
\item For odd $m$, use $\sin^2(x)=1-\cos^2(x)$ and then substitute $u=\cos(x),\quad\dd{u}=-\sin(x)\dd{x}$.
\item For $m=n$, use $\sin(x)\cos(x)=\frac{1}{2}\sin(2x)$.
\item For $\frac{m-n}{2}\in\bbN$, use $\sin(x)\cos(x)=\frac{1}{2}\sin(2x)$ and then substitute $\sin^2(x)=\frac{1}{2}(1-\cos(2x))$.
\item For $\frac{n-m}{2}\in\bbN$, use $\sin(x)\cos(x)=\frac{1}{2}\sin(2x)$ and then substitute $\cos^2(x)=\frac{1}{2}(1+\cos(2x))$.
\end{itemize}
\subsubsection{Product of tangent and secant functions}
\[\int\tan^m(x)\sec^n(x)\dd{x},\quad m,n\in\bbN.\]
\begin{itemize}
\item For even $n$, save a factor of $\sec^2(x)$, use $\sec^2(x)=1+\tan^2(x)$ to express the remaining factors in terms of $\tan(x)$, and then substitute $u=\tan(x),\quad\dd{u}=\sec^2(x)\dd{x}$.
\item For odd $m$, save a factor of $\tan(x)\sec(x)$, use $\tan^2(x)=\sec^2(x)-1$ to express the remaining factors in terms of $\sec(x)$, and then substitute $u=\sec(x),\quad\dd{u}=\tan(x)\sec(x)\dd{x}$.
\item For even $m$, use $\tan^2(x)=\sec^2(x)-1$ to express the integrand in terms of $\sec(x)$, and then use reduction formula for secant function.
\end{itemize}
\subsubsection{Product of cotangent and cosecant functions}
\[\int\cot^m(x)\csc^n(x)\dd{x},\quad m,n\in\bbN.\]
\begin{itemize}
\item For even $n$, save a factor of $\csc^2(x)$, use $\csc^2(x)=1+\cot^2(x)$ to express the remaining factors in terms of $\cot(x)$, and then substitute $u=\cot(x),\quad\dd{u}=-\csc^2(x)\dd{x}$.
\item For odd $m$, save a factor of $\cot(x)\csc(x)$, use $\cot^2(x)=\csc^2(x)-1$ to express the remaining factors in terms of $\csc(x)$, and then substitute $u=\csc(x),\quad\dd{u}=-\cot(x)\csc(x)\dd{x}$.
\item For even $m$, use $\cot^2(x)=\csc^2(x)-1$ to express the integrand in terms of $\csc(x)$, and then use reduction formula for cosecant function.
\end{itemize}
\subsection{Leibniz integral rule (for differentiation under the integral sign) (（積分符號內取微分的）萊布尼茲積分法則)}
\subsubsection{Constant limits}
Let $f(x,t)$ be a function of two variables and the following expression exists. Then:
\[\dv{}{x}\qty(\int_a^bf(x,t)\dd{t})=\int_a^b\pdv{}{x}f(x,t)\dd{t}.\]
\subsubsection{Variable limits}
Let $a(x),b(x)$ be two functions of one variable, $f(x,t)$ be a function of two variables, and the following expression exists. Then:
\[\dv{}{x}\qty(\int_{a(x)}^{b(x)}f(x,t)\dd{t})=f\qty(x,b(x))\cdot\dv{b(x)}{x}-f\qty(x,a(x))\cdot\dv{a(x)}{x}+\int_{a(x)}^{b(x)}\pdv{}{x}f(x,t)\dd{t}\]
\subsubsection{Measure theory form for const limits}
Let $X$ be an open subset of $\mathbb{R}$, and $\Omega$ be a measure space. Suppose $f\colon X\times\Omega\to\mathbb{R}$ satisfies the following conditions:
\ben
\item $f(x,\omega)$ is a Lebesgue-integrable function of $\omega$ for each $x\in X$.
\item For almost all $\omega\in\Omega$, the partial derivative $\pdv{f}{x}$ exists for all $x\in X$.
\item There is an integrable function $\theta\colon\Omega\to\mathbb{R}$ such that $\abs{\pdv{f}{x}\qty(x,\omega)}\leq\theta(\omega)$ for all $x\in X$ and almost all $\omega\in\Omega$.
\een
Then, for all $x\in X$,
\[\dv{}{x}\int_{\Omega}f(x,\omega)\dd{\omega}=\int_{\Omega}\pdv{f}{x}\qty(x,\omega)\dd{\omega}.\]
\ssc{Lebesgue integral theorems}
\sssc{Fubini's theorem (富比尼定理)}
If real-valued function $f\in L^1(X\times Y)$, then:
\[\int_{X\times Y}f(x,y)\dd{(x,y)}=\int_X\int_Yf(x,y)\dd{y}\dd{x}=\int_Y\int_Xf(x,y)\dd{x}\dd{y}.\]
\bpr
For nonnegative simple function, it is obvious.

For a nonnegative measurable function $f$, construct a monotone series of nonnegative simple functions $f_n$ that converges pointwise to $f$:
\[f_n=\sum_{k=0}^{2^nn-1}\frac{k}{2^n}\cdot\mb{1}_{\left\{x\mid\frac{x}{2^n}\le f(x)<\frac{k+1}{2^n}\right\}}(x)+n\cdot 1_{\{x\mid f(x)\ge n\}}(x)\]
with
\[\lim_{n\to\infty}f_n=f.\]
By monotone convergence theorem, the statement holds for all nonnegative measurable functions.

For measurable function, define
\[\begin{aligned}
f^{+}(x)&=
\begin{cases}
f(x),\quad&\text{if\ }f(x)>0\\
0,\quad &\text{otherwise}
\end{cases},\quad\tx{and}\\
f^{-}(x)&=
\begin{cases}
-f(x),\quad&\text{if\ }f(x)<0\\
0,\quad&\text{otherwise}
\end{cases}.
\end{aligned}\]
Note that both $f^+$ and $f^-$ are nonnegative and that
\[f=f^+-f^-,\quad |f|=f^++f^-.\]
Thus, the statement holds for all measurable functions.
\epr
\sssc{Countable additivity}
For an indexed family $(A_i)_{i=1}^n$ of piecewise disjoint measurable sets with $n\in\bbN\cup\{\infty\}$, if $f\in L^1(A_i)$ for all $i$, then
\[\int_{\bigcup_{i=1}^nA_i}f\dd{\mu}=\sum_{i=1}^n\int_{A_i}f\dd{\mu}.\]
\sssc{(Lebesgue's) Dominated convergence theorem (DCT) (（勒貝格）控制/受制收斂定理)}
Let $\langle f_n\rangle_{n\in I}$ be a net of real or complex-valued measurable functions on a measure space $(S,\Sigma,\mu)$. If $f_n$ is almost everywhere pointwise convergent to a function $f$, and there is a Lebesgue-integrable function $g$ such that
\[\abs{f_n}\leq g\]
almost everywhere for all $n\in I$.

Then $f_n$ for all $n\in I$ and $f$ are in $L^1(\mu)$ and
\[\lim_n\int_Sf_n\dd{\mu}=\int_S\lim_nf_n\dd{\mu}=\int_Sf\dd{\mu},\]
and
\[\lim_n\int_S\abs{f_n-f}\dd{\mu}=0.\]
\bpr
PLACEHOLDER
\epr
\subsection{Fundamental theorem of multivariable calculus (多變數微積分基本定理)}
\subsubsection{Gradient theorem or Fundamental theorem for linear integral}
Suppose $C$ is a piecewise $C^1$, oriented curve with initial point $a$ and terminal point $b$. If $\mb{F}$ is a scalar or vector field differentiable on a neighborhood of $C$, then
\[\int_C(\nabla\mathbf{F})\cdot\dd{\mb{r}}=\mb{F}(b)-\mb{F}(a).\]
\bpr
Let $C$ be $\mb{r}(t),\quad t\in[\alpha,\beta]$, where $\mb{r}(\alpha)=a$ and $\mb{r}(\beta)=b$.
\[\int_C(\nabla\mathbf{F})\cdot\dd{\mb{r}}=\int_{\alpha}^{\beta}(\nabla\mb{F})(\mb{r}(t))\mb{r}'(t)\dd{t}\]
Let $\mb{F}=(F_1,\ldots,F_n)$. For each component $F_i(\mb{r}(t))$,
\[\dv{}{t}F_i(\mb{r}(t))=\nabla F_i(\mb{r}(t))\cdot\mb{r}'(t)\]
Thus
\[(\nabla\mb{F})(\mb{r}(t))\mb{r}'(t)=\dv{}{t}F_i(\mb{r}(t))\]
By FTC part 2,
\[\int_{\alpha}^{\beta}\dv{}{t}\mb{F}(\mb{r}(t))\dd{t}=\mb{F}(\mb{r}(\beta))-\mb{F}(\mb{r}(\alpha))=\mb{F}(b)-\mb{F}(a)\]
\epr
\subsubsection{Divergence theorem, Gauss's theorem, or Ostrogradsky's theorem (高斯散度定理)}
Suppose $V\subseteq\mathbb{R}^n$ is a compact volume with outward-oriented, piecewise $C^1$ boundary $\partial V=S$. If $\mb{F}$ is a vector field differentiable and piecewise $C^1$ on a neighborhood of $V$, then
\[\int_V\nabla\cdot\mb{F}\dd{V}=\oint_S\mb{F}\cdot\dd{\mb{S}}.\]
\bpr
Simple regions case: Let $\sigma$ be a permutation on $(x_1,\ldots,x_n)$ and
\[V=\left\{(x_1,\ldots,x_n)\colon\forall i\in\bbZ\cap[1,n]\colon\varphi_i\qty(x_{\sigma(j)}_{j=1}^i)\le x_{\sigma(i)}\le\psi_i\qty(x_{\sigma(j)}_{j=1}^i)\right\},\]
where $\varphi_i,\psi_i$ are $(i-1)$-ary $C^1$ functions.
\[\int_V\nabla\cdot\mb{F}\dd{V}=\sum_{i=1}^n\int_V\pdv{F_{\sigma(i)}}{x_{\sigma(i)}}\dd{V}\]
For each term, by FTC part 2,
\[\ba
\int_V\pdv{F_i}{x_i}\dd{V}=&\int_{\varphi_1}^{\psi_1}\ldots\int_{\varphi_n(x_{\sigma(1)},\ldots,x_{\sigma(n-1)})}^{\psi_n(x_{\sigma(1)},\ldots,x_{\sigma(n-1)})}\pdv{F_{\sigma(i)}}{x_{\sigma(i)}}\dd{x_{\sigma(n)}}\ldots\dd{x_{\sigma(1)}}\\
=&\int_{\varphi_1}^{\psi_1}\ldots\int_{\varphi_{i-1}(x_{\sigma(1)},\ldots,x_{\sigma(i-2)})}^{\psi_{i-1}(x_{\sigma(1)},\ldots,x_{\sigma(i-2)})}\int_{\varphi_{i+1}(x_{\sigma(1)},\ldots,x_{\sigma(i)})}^{\psi_{i+1}(x_{\sigma(1)},\ldots,x_{\sigma(i)})}\ldots\int_{\varphi_n(x_{\sigma(1)},\ldots,x_{\sigma(n-1)})}^{\psi_n(x_{\sigma(1)},\ldots,x_{\sigma(n-1)})}F_{\sigma(i)}}\qty(x_1,\ldots,x_{i-1},\psi_i(x_{\sigma(1)},\ldots,x_{\sigma(i-1)}),x_{i+1},\ldots,x_n)-F_{\sigma(i)}}\qty(x_1,\ldots,x_{i-1},\varphi_i(x_{\sigma(1)},\ldots,x_{\sigma(i-1)}),x_{i+1},\ldots,x_n)\dd{x_{\sigma(n)}}\ldots\dd{x_{\sigma(i+1)}}\dd{x_{\sigma(i-1)}}\ldots\dd{x_{\sigma(1)}}\\
=&\int_{\varphi_1}^{\psi_1}\ldots\int_{\varphi_{i-1}(x_{\sigma(1)},\ldots,x_{\sigma(i-2)})}^{\psi_{i-1}(x_{\sigma(1)},\ldots,x_{\sigma(i-2)})}\int_{\varphi_{i+1}(x_{\sigma(1)},\ldots,x_{\sigma(i)})}^{\psi_{i+1}(x_{\sigma(1)},\ldots,x_{\sigma(i)})}\ldots\int_{\varphi_n(x_{\sigma(1)},\ldots,x_{\sigma(n-1)})}^{\psi_n(x_{\sigma(1)},\ldots,x_{\sigma(n-1)})}
(\mb{F}\cdot\mb{e}_i)\qty(\psi_i(x_{\sigma(1)},\ldots,x_{\sigma(i-1)}))+(\mb{F}\cdot(-\mb{e}_i))\qty(\varphi_i(x_{\sigma(1)},\ldots,x_{\sigma(i-1)}))\dd{x_{\sigma(n)}}\ldots\dd{x_{\sigma(i+1)}}\dd{x_{\sigma(i-1)}}\ldots\dd{x_{\sigma(1)}}\\
=&\int_{x_{\sigma(i)}=\psi_i(x_{\sigma(1)},\ldots,x_{\sigma(i-1)})}\mb{F}\cdot\dd{\mb{S}}+\int_{x_i=\varphi_i(x_{\sigma(1)},\ldots,x_{\sigma(i-1)})}\mb{F}\cdot\dd{\mb{S}}
\ea\]
Thus
\[\sum_{i=1}^n\int_V\pdv{F_i}{x_i}\dd{V}=\int_S\mb{F}\cdot\dd{\mb{S}}\]
This proves for simple regions case.

Finite union of simple regions case: Both sides are additive. Internal faces cancel. So the identity still holds.

Any compact region with piecewise $C^1$ boundary has finite cover of simple regions. So the identity still holds.
\epr
\subsubsection{Green's theorem (格林定理或綠定理)}
Suppose $S\subseteq\bbR^2$ is a compact surface with positively-oriented, piecewise $C^1$ boundary $\partial S=\ell$. If two scalar functions $P(x,y)$ and $Q(x,y)$ are continuous and piecewise $C^1$ on a neighborhood of $S$, then
\[\iint_S\pdv{Q}{x}-\pdv{P}{y}\dd{S}=\oint_{\ell}\qty(P\dd{x}+Q\dd{y}).\]
\bpr
Let $\mb{F}=(Q,-P)$.
\[\nabla\cdot\mb{F}=\pdv{Q}{x}-\pdv{P}{y}\]
\[\dd{\boldsymbol{\ell}}=\dd{y}-\dd{x}\]
By divergence theorem,
\[\iint_S\nabla\cdot\mb{F}\dd{S}=\oint_{\ell}\mb{F}\cdot\dd{\boldsymbol{\ell}}.\]
\epr
\subsubsection{Stokes' theorem (斯托克斯定理) or Kelvin–Stokes theorem}
Suppose $S\subseteq\bbR^3$ is a compact surface with positively-oriented, piecewise $C^1$ boundary $\partial S=\ell$. If a three-dimensional vector field $\mb{F}$ is continuous and piecewise $C^1$ on a neighborhood of $S$, then
\[\iint_S(\nabla\times\mathbf{F})\cdot \dd{\mathbf{S}}=\oint_{\ell}\mathbf{F}\cdot\dd{\boldsymbol{\ell}}.\]
\bpr
Parameteriz $S$ and $\ell$ as $T$ and $\omega$ in parameters $u,v$. Let $T$ be $\mb{r}(u,v)=(x(u,v),y(u,v),z(u,v))$, and
\[\dd{\mb{T}}=\pdv{\mb{r}}{u}\times\pdv{\mb{r}}{v}\dd{u}\dd{v}.\]
\[\iint_S(\nabla\times\mb{F})\dd{\mb{S}}=\iint_T(\nabla\times\mb{F})(\mb{r}(u,v))\cdot\qty(\pdv{\mb{r}}{u}\times\pdv{\mb{r}}{v})\dd{u}\dd{v}\]
Define the pulled-back vector fields in parameter space
\[\mb{G}(u,v)=\mb{F}(\mb{r}(u,v))\]
\[\mb{P}(u,v)=\mb{F}\cdot\pdv{\mb{r}}{u}\]
\[\mb{Q}(u,v)=\mb{F}\cdot\pdv{\mb{r}}{v}\]
Let $\mb{F}=(F_1,F_2,F_3)$.
\[\ba
&{}(\nabla\times\mb{F})(\mb{r}(u,v))\cdot\qty(\pdv{\mb{r}}{u}\times\pdv{\mb{r}}{v})\\
=&\begin{vmatrix}
\pdv{F_3}{y}-\pdv{F_2}{z} & \pdv{F_1}{z}-\pdv{F_3}{x} & \pdv{F_2}{x}-\pdv{F_1}{y}\\
\pdv{x}{u} & \pdv{y}{u} & \pdv{z}{u}\\
\pdv{x}{v} & \pdv{y}{v} & \pdv{z}{v}
\end{vmatrix}\\
=&\qty(\pdv{F_3}{y}-\pdv{F_2}{z})\qty(\pdv{y}{u}\pdv{z}{v}-\pdv{y}{v}\pdv{z}{u})-\qty(\pdv{F_1}{z}-\pdv{F_3}{x})\qty(\pdv{x}{u}\pdv{z}{v}-\pdv{x}{v}\pdv{z}{u})+\qty(\pdv{F_2}{x}-\pdv{F_1}{y})\qty(\pdv{x}{u}\pdv{y}{v}-\pdv{x}{v}\pdv{y}{u})\\
=&\pdv{F_3}{y}\pdv{y}{u}\pdv{z}{v}-\pdv{F_3}{y}\pdv{y}{v}\pdv{z}{u}-\pdv{F_2}{z}\pdv{y}{u}\pdv{z}{v}+\pdv{F_2}{z}\pdv{y}{v}\pdv{z}{u}-\pdv{F_1}{z}\pdv{x}{u}\pdv{z}{v}+\pdv{F_1}{z}\pdv{x}{v}\pdv{z}{u}-\pdv{F_3}{x}\pdv{x}{u}\pdv{z}{v}+\pdv{F_3}{x}\pdv{x}{v}\pdv{z}{u}+\pdv{F_2}{x}\pdv{x}{u}\pdv{y}{v}-\pdv{F_2}{x}\pdv{x}{v}\pdv{y}{u}-\pdv{F_1}{y}\pdv{x}{u}\pdv{y}{v}+\pdv{F_1}{y}\pdv{x}{v}\pdv{y}{u}\\
=&\qty(\pdv{F_1}{x}\pdv{x}{u}+\pdv{F_1}{y}\pdv{y}{u}+\pdv{F_1}{z}\pdv{z}{u})\pdv{x}{v}-\pdv{F_1}{x}\pdv{x}{u}\pdv{x}{v}-\qty(\pdv{F_1}{x}\pdv{x}{v}+\pdv{F_1}{y}\pdv{y}{v}+\pdv{F_1}{z}\pdv{z}{v})\pdv{x}{u}+\pdv{F_1}{x}\pdv{x}{u}\pdv{x}{v}+\qty(\pdv{F_2}{x}\pdv{x}{u}+\pdv{F_2}{y}\pdv{y}{u}+\pdv{F_2}{z}\pdv{z}{u})\pdv{x}{v}-\pdv{F_2}{y}\pdv{y}{u}\pdv{y}{v}-\qty(\pdv{F_2}{x}\pdv{x}{v}+\pdv{F_2}{y}\pdv{y}{v}+\pdv{F_2}{z}\pdv{z}{v})\pdv{x}{u}+\pdv{F_2}{y}\pdv{y}{u}\pdv{y}{v}+\qty(\pdv{F_3}{x}\pdv{x}{u}+\pdv{F_3}{y}\pdv{y}{u}+\pdv{F_3}{z}\pdv{z}{u})\pdv{x}{v}-\pdv{F_3}{z}\pdv{z}{u}\pdv{z}{v}-\qty(\pdv{F_3}{x}\pdv{x}{v}+\pdv{F_3}{y}\pdv{y}{v}+\pdv{F_3}{z}\pdv{z}{v})\pdv{x}{u}+\pdv{F_3}{z}\pdv{z}{u}\pdv{z}{v}\\
=&\qty(\pdv{F_1}{u}\pdv{x}{v}+\pdv{F_2}{u}\pdv{y}{v}+\pdv{F_3}{u}\pdv{z}{v})-\qty(\pdv{F_1}{v}\pdv{x}{u}+\pdv{F_2}{v}\pdv{y}{u}+\pdv{F_3}{v}\pdv{z}{u})\\
=&\pdv{}{u}\qty(\mb{F}\cdot\pdv{\mb{r}}{v})-\pdv{}{v}\qty(\mb{F}\cdot\pdv{\mb{r}}{u})\\
=&\pdv{\mb{Q}}{u}-\pdv{\mb{P}}{v}
\ea\]
By Green's theorem,
\[\ba
\iint_T\pdv{\mb{Q}}{u}-\pdv{\mb{P}}{v}\dd{u}\dd{v}=&\oint_{\omega}\qty(\mb{P}\dd{u}+\mb{Q}\dd{v})\\
=&\oint_{\omega}\mb{G}\dd{\boldsymbol{\omega}}\\
=&\oint_{\ell}\mb{F}\dd{\boldsymbol{\ell}}
\ea\]
\epr
\subsubsection{Generalized Stokes theorem, Stokes–Cartan theorem, or fundamental theorem of multivariable calculus}
The generalized Stokes theorem says that the integral of a differential form $\omega$ over the boundary $\partial\Omega$ of some orientable manifold $\Omega$ is equal to the integral of its exterior derivative $\mathrm{d}\boldsymbol{\omega}$ over the whole of $\Omega$, i.e.,
\[\int _{\partial\Omega}\omega=\int_{\Omega}\mathrm{d}\boldsymbol{\omega}\]
\bpr
PLACEHOLDER
also statement
\epr
\ssc{Conservative field}
\sssc{Path independence equivalent to zero closed path integral}
For an $n$-dimensional vector field $\mb{v}$ defined on an open subset $U$ of $\bbR^n$, for any two piecewise $C^1$ paths $C,D\subseteq U$ with same initial and terminal points,
\[\int_C\mb{v}\cdot\dd{\mb{x}}=\int_D\mb{v}\cdot\dd{\mb{x}}\]
iff for any piecewise $C^1$ closed path $E\subseteq U$,
\[\oint_E\mb{v}\cdot\dd{\mb{x}}=0.\]
\bpr
By $\int_C\mb{v}\cdot\dd{\mb{x}}=-\int_{-C}\mb{v}\cdot\dd{\mb{x}}$ and $\int_{C+D}\mb{v}\cdot\dd{\mb{x}}=\int_{C}\mb{v}\cdot\dd{\mb{x}}+\int_{D}\mb{v}\cdot\dd{\mb{x}}$.
\epr
\sssc{Conservative equivalent to path independence}
For a conservative field $f$ on an open subset $U$ of $\bbR^n$ and any two piecewise $C^1$ paths $C,D\subseteq U$ with same initial point $a$ and terminal point $b$,
\[\int_Cf\cdot\dd{\mb{x}}=\int_Df\cdot\dd{\mb{x}}.\]
\bpr
Let $f=\nabla\varphi$, by Gradient theorem, both equal to $\varphi(b)-\varphi(a)$.
\epr
\sssc{One dimension}
A one-dimensional field is conservative on an open interval iff it is continuous on it.
\sssc{Two dimensions}
If a two-dimensional field $f=(f_1,f_2)$ is conservative on a connected open set $U$, then
\[\pdv{f_2}{x}-\pdv{f_1}{y}=0\]
on the subset of $U$ where LHS exists.
\bpr
Let $f=\nabla\varphi$. $f_1=\varphi_x$, $f_2=\varphi_y$,
\[\pdv{f_2}{x}=\pdv{f_1}{y}=\varphi_{xy}.\]
\epr
A $C^1$ two-dimensional field $f=(f_1,f_2)$ is conservative on a simply-connected open set $U$ iff
\[\pdv{f_2}{x}-\pdv{f_1}{y}=0\]
on $U$.
\bpr
If: By Green's theorem, let $C$ be a piecewise smooth curve in $U$ and surrounds surface $S$.
\[\oint_Cf\cdot\dd{(x,y)}=\iint_S\pdv{f_2}{x}-\pdv{f_1}{y}\dd{S}=0\]
\epr
\sssc{Three dimensions}
If a three-dimensional field $f$ is conservative on a connected open set $U$, then
\[\nabla\times f=\mb{0}\]
on the subset of $U$ where LHS exists.
\bpr
Let $f=\nabla\varphi=(f_1,f_2,f_3)$. $f_1=\varphi_x$, $f_2=\varphi_y$, $f_3=\varphi_z$,
\[\ba
\nabla\times f=&\qty(\pdv{f_3}{y}-\pdv{f_2}{z},\pdv{f_1}{z}-\pdv{f_3}{x},\pdv{f_2}{x}-\pdv{f_1}{y})\\
=&\qty(\varphi_{yz}-\varphi_{yz},\varphi_{xz}-\varphi_{xz},\varphi_{xy}-\varphi_{xy})\\
=&(0,0,0)
\ea\]
\epr
A $C^1$ three-dimensional field $f$ is conservative on a simply-connected open set $U$ iff
\[\nabla\times f=\mb{0}\]
on $U$.
\bpr
If: By Stokes' theorem, let $C$ be a piecewise smooth curve in $U$ and surrounds surface $S$.
\[\oint_Cf\cdot\dd{(x,y,z)}=\iint_S\nabla\times f\dd{S}=0\]
\epr
\ssc{Divergence of curl}
For a two times differentiable three-dimensional field $\mb{T}$, $\nabla\cdot(\nabla\times\mb{T})=0.\]
\bpr
\[\ba
\nabla\cdot(\nabla\times\mb{T})=&\nabla\cdot\qty(\pdv{T_3}{y}-\pdv{T_2}{z},\pdv{T_1}{z}-\pdv{T_3}{x},\pdv{T_2}{x}-\pdv{T_1}{y})\\
=&\pdv[2]{T_3}{xy}-\pdv[2]{T_2}{xz}+\pdv{T_1}{yz}-\pdv{T_3}{xy}+\pdv{T_2}{xz}-\pdv{T_1}{yz}\\
=&0
\ea\]
\epr
\ssc{Two-dimensional curve rotates about principal axis on the same plane}
The area of the surface formed by rotating the curve $\mb{r}(t)=(x(t),y(t)),\quad t\in[a,b]$ about $y=d$ is
\[A=2\pi\int_{t=a}^b\abs{y(t)-d}\dd{s(t)},\]
where $s(t)$ is the arc length function of $\mb{r}$ from $a$ to $t$.
\subsection{Solid of revolution}
\subsubsection{Pappus's (centroid) theorem I or (Pappus–)Guldinus theorem I}
The area of the surface formed by rotating a two-dimensional curve with arc length $s$ and centroid $\mb{P}$ about a rotational axis on the same plane of the curve that doesn't cut the curve in a distance $d$ from $\mb{P}$ is
\[A=2\pi sd.\]
\begin{proof}
Let the curve be $\mb{r}(t)=(x(t),y(t)),\quad t\in[a,b]$, $s(t)$ be the arc length function of it, the centroid of it be the origin, that is, 
\[\int\mb{r}(t)\dd{s(t)}=(0,0),\]
and the rotational axis is $y=d$.
\[\ba
A&=2\pi\int_{t=a}^b\abs{y(t)-d}\dd{s(t)}\\
&=2\pi\abs{\int_{t=a}^by(t)\dd{s(t)}-\int_{t=a}^bd\dd{s(t)}}\\
&=2\pi sd
\ea\]
\end{proof}
\subsubsection{Disc method or disc integration}
The volume of the solid formed by rotating the area bounded by $y=f(x)$,$y=g(x)$, $x=a$, and $x=b$ with $a<b$ about the $x$-axis is given by 
\[V=\pi\int_a^b\abs{f(x)^2-g(x)^2}\dd{x},\]
aka Washer method.

If $g(x)=0$, this reduces to: 
\[V=\pi\int_a^bf(x)^2\dd{x}.\]
\subsubsection{Shell method or shell integration}
The volume of the solid formed by rotating the area bounded by $y=f(x)$,$y=g(x)$, $x=a$, and $x=b$ with $a<b$ about the $y$-axis is given by 
\[V=2\pi\int_a^bx\abs{f(x)-g(x)}\dd{x}.\]
\subsubsection{Pappus's (centroid) theorem II or (Pappus–)Guldinus theorem II}
The volume of the solid formed by rotating a plane region with area $A$ and centroid $\mb{P}$ about a rotational axis on the same plane of the region that doesn't cut the region in a distance $d$ from $\mb{P}$ is
\[V=2\pi Ad.\]
\begin{proof}
Let the rotational axis be the $y$-axis. Split the region to subregions such that each subregion is bounded by $y=f(x)$,$y=g(x)$, $x=a$, and $x=b$, for some functions $f(x)$ and $g(x)$ and some constant $a<b$.

For each subregion,
\[V=2\pi\int_a^bx\abs{f(x)-g(x)}\dd{x}.\]
\[Ad=\int_a^bx\abs{f(x)-g(x)}\dd{x}.\]
\[V=2\pi Ad.\]

The sum of the volume of the solid formed by rotating the subregions is the volume of the solid formed by rotating the original region. The average of the centroids $\mb{P}_i=(d_i,y_i)$ of the subregions weighted by their areas $A_i$ is the centroid $\mb{P}=(d,p_y)$ of the original region.
\[\mb{P}=\frac{\sum_iA_i\mb{P}_i}{A}\]
\[V=\sum_iV_i=\sum_iA_id_i=A\frac{\sum_iA_id_i}{A}=Ad.\]
\end{proof}
\subsection{Numerical integration (數值積分)}
\subsubsection{Left and right endpoint rule}
Let $f$ be a continuous real-valued function on $[a,b]$. The left endpoint rule gives the approximation of $\int_a^bf(x)\dd{x}$
\[L_n=\frac{b-a}{n}\sum_{i=0}^{n-1}f\qty(a+\frac{i}{n}(b-a)).\]
The right endpoint rule gives the approximation of $\int_a^bf(x)\dd{x}$
\[R_n=\frac{b-a}{n}\sum_{i=1}^nf\qty(a+\frac{i}{n}(b-a)).\]
\subsubsection{Midpoint rule}
Let $f$ be a continuous real-valued function on $[a,b]$. The midpoint rule gives the approximation of $\int_a^bf(x)\dd{x}$
\[M_n=\frac{b-a}{n}\sum_{i=0}^{n-1}f\qty(a+\frac{2i+1}{2n}(b-a)).\]
The error is
\[E_M=\int_a^bf(x)\dd{x}-M_n.\]
If $f''$ exists a.e. on $[a,b]$ and $\abs{f''(x)}\leq K$ for all $x\in[a,b]$ such that $f''(x)$ exist, then
\[\abs{E_M}\leq\frac{K(b-a)^3}{24n^2}.\]
If $f\in C^2([a,b])$, then
\[E_M=\frac{(b-a)^3}{24n^2}f''(\xi)\]
for some $\xi\in[a,b]$.
\subsubsection{Trapezoidal rule}
Let $f$ be a continuous real-valued function on $[a,b]$. The trapezoidal rule gives the approximation of $\int_a^bf(x)\dd{x}$
\[T_n=\frac{b-a}{2n}\qty(f(a)+f(b)+2\sum_{i=1}^{n-1}f\qty(a+\frac{i}{n}(b-a))).\]
The error is
\[E_T=\int_a^bf(x)\dd{x}-T_n.\]
If $f''$ exists a.e. on $[a,b]$ and $\abs{f''(x)}\leq K$ for all $x\in[a,b]$ such that $f''(x)$ exist, then
\[\abs{E_T}\leq\frac{K(b-a)^3}{12n^2}.\]
If $f\in C^2([a,b])$, then
\[E_T=-\frac{(b-a)^3}{12n^2}f''(\xi)\]
for some $\xi\in[a,b]$.
\subsubsection{Simpson's rule or Simpson's 1/3 rule}
Let $f$ be a continuous real-valued function on $[a,b]$, the Simpson's rule gives the approximation of $\int_a^bf(x)\dd{x}$
\[S_n=\frac{b-a}{3n}\qty(f(a)+f(b)+2\sum_{i=1}^{n-1}f\qty(a+\frac{i(i\mod 2+1)}{n}(b-a)),\quad\frac{n}{2}\in\bbN.\]
The error is
\[E_S=\int_a^bf(x)\dd{x}-S_n.\]
If $f^{(4)}$ exists a.e. on $[a,b]$ and $\abs{f^{(4)}(x)}\leq K$ for all $x\in[a,b]$ such that $f^{(4)}(x)$ exist, then
\[\abs{E_S}\leq\frac{K(b-a)^5}{180n^4}.\]
If $f\in C^4([a,b])$, then
\[E_S=-\frac{(b-a)^5}{180n^4}f^{(4)}(\xi)\]
for some $\xi\in[a,b]$.
\ssc{Symmetry function integral}
Suppose $f\colon\bbR\to\bbR$ is such that $f(x)=f(-x)$ and $f\in L^1(\bbR)$. Let
\[I=\int_{-\infty}^{\infty}f(x)\dd{x}\]
and
\[F(x)=\int_{-\infty}^xf(t)\dd{t}\]
Then,
\[F(x)+F(-x)=I\]
\bpr
For $x\ge 0$,
\[F(x)=\frac{I}{2}+\int_0^xf(t)\dd{t}\]
\[F(-x)=\int_{-\infty}^{-x}f(t)\dd{t}=\int_x^{\infty}f(t)\dd{t}\]
\epr
\ssc{Multiple integral theorems}
\sssc{Multiple integral to multiplication of integrals}
For a rectangle
\[R=\{(x_1,x_2,\ldots,x_n)\mid a_1\le x_1\le b_1\land a_2\le x_2\le b_2\land\ldots a_n\le x_n\le b_n\},\]
\[\int_R\prod_{i=1}^nf_i(x_i)\dd{V}=\prod_{i=1}^n\int_{a_i}^{b_i}f_i(x)\dd{x}.\]
\sssc{Multiple integral to iterated integral theorem}
For a region $D$ in $\bbR^n$ of the form
\[D=\left\{(x_1,x_2,\ldots,x_n)\middle|\bigwedge_{i=0}^{n-1}f_i\qty(x_{\sigma(1)},x_{\sigma(2)},\ldots,x_{\sigma(i)})\le x_{\sigma(i+1)}\le g_i\qty(x_{\sigma(1)},x_{\sigma(2)},\ldots,x_{\sigma(i)}),\]
where $\sigma$ is any permutation on $[1,n]\cap\bbN$, $f_0,g_0$ are real numbers, and $f_i,g_i$ are $i$-ary functions continuous on $\qty[f_{i-1}\qty(x_{\sigma(1)},x_{\sigma(2)},\ldots,x_{\sigma(i-1)}),g_{i-1}\qty(x_{\sigma(1)},x_{\sigma(2)},\ldots,x_{\sigma(i-1)})]$, and function $F\in L^1(D)$,
\[\int_DF(x_1,x_2,\ldots,x_n)\dd{\mb{x}}=\qty(\int_{f_i\qty(x_{\sigma(1)},x_{\sigma(2)},\ldots,x_{\sigma(i)})}^{g_i\qty(x_{\sigma(1)},x_{\sigma(2)},\ldots,x_{\sigma(i)})})_{i=0}^{n-1}F(x_1,x_2,\ldots,x_n)\qty(\dd{x_{\sigma(n-i)})_{i=0}^{n-1}.\]
